% 第4章：整数阈值的排序扫描。
% 样本（已按输入排序并合并重复）：x=2:负1；x=5:正1、负1；x=6:正1；x=9:正1；x=11:负1。
% 切分 c_k 位于第 k 个与第 k+1 个不同输入之间，左侧判 0、右侧判 1。
% 错误数 = 左侧正例数 + 右侧负例数：c_0..c_5 依次为 3, 2, 2, 1, 0+... 见下表：
%   c0 (t<=2): 左正0 + 右负3 = 3
%   c1 (2<t<=5): 0 + 2 = 2
%   c2 (5<t<=6): 1 + 1 = 2
%   c3 (6<t<=9): 2 + 1 = 3
%   c4 (9<t<=11): 3 + 1 = 4
%   c5 (t>11): 3 + 0 = 3
% 最少错误 2，在 c1 与 c2 并列；选较小阈值得到 t = 3（与 t=3,4,5 等价）。
% 越过负例错误数减1，越过正例加1；x=5 同时含正、负，越过它净变化为 0。
\begingroup
\definecolor{FigureBlue}{HTML}{7188A3}\definecolor{FigureBlueFill}{HTML}{E5EAF0}
\definecolor{FigureGreen}{HTML}{829C8B}\definecolor{FigureGreenFill}{HTML}{E8EEE9}
\definecolor{FigureYellow}{HTML}{B4A16F}\definecolor{FigureYellowFill}{HTML}{F1ECDD}
\definecolor{FigureRed}{HTML}{A7787D}\definecolor{FigureRedFill}{HTML}{F0E5E6}
\definecolor{FigureInk}{HTML}{3B4252}\definecolor{FigureMuted}{HTML}{6B7280}
\begin{tikzpicture}[figure picture,x=1mm,y=1mm]
\def\u{7.5}
\draw[draw=FigureStroke,line width=.75pt,-{Stealth[length=1.6mm]}] (\u*0.5,10)--(\u*12.6,10);
\foreach \v in {1,...,12}{\draw[draw=FigureStroke] (\u*\v,9.3)--(\u*\v,10.7);\node[anchor=north,font=\figurefont\scriptsize,text=FigureMuted] at (\u*\v,9) {$\v$};}
\node[anchor=west] at (\u*12.6,10) {$x$};
% observations: negatives hollow blue, positives filled red, stacked
\foreach \x/\yy in {2/14,5/14,11/14}{\draw[draw=FigureBlue,fill=FigureBlueFill,line width=.7pt] (\u*\x,\yy) circle(1.3);}
\foreach \x/\yy in {5/17.4,6/14,9/14}{\draw[draw=FigureRed,fill=FigureRed,line width=.7pt] (\u*\x,\yy) circle(1.3);}
% cut positions and error counts
\foreach \c/\xc/\e in {0/1.2/3,1/3.5/2,2/5.5/2,3/7.5/3,4/10/4,5/11.6/3}{
  \draw[draw=FigureMuted,line width=.5pt,dash pattern=on 1pt off 1pt] (\u*\xc,11.4)--(\u*\xc,21);
  \node[font=\figurefont\scriptsize] at (\u*\xc,24) {$c_{\c}$};
  \node at (\u*\xc,1.8) {$\e$};}
\fill[FigureYellowFill] (\u*3.5-3,-0.6) rectangle (\u*5.5+3,4.2);
\node at (\u*3.5,1.8) {$2$};\node at (\u*5.5,1.8) {$2$};
\node[anchor=east] at (\u*0.9,1.8) {错误数};
\node[figure note,anchor=west] at (\u*12.6,19) {空心：负例\\实心：正例};
\end{tikzpicture}
\endgroup
